\documentclass[aspectratio=169]{beamer}

\usetheme{default}
\setbeamertemplate{navigation symbols}{}
\setbeamercolor{background canvas}{bg=white}
\setbeamertemplate{footline}{}
\setbeamertemplate{headline}{}

\usepackage[T1]{fontenc}
\usepackage{tikz}
\usetikzlibrary{calc}
\usepackage{graphicx}
\usepackage{tcolorbox}
\usepackage{ragged2e}

\definecolor{hciblue}{HTML}{1F618D}
\definecolor{cgred}{HTML}{C0392B}
\definecolor{darktext}{HTML}{1A1A1A}
\definecolor{graytext}{HTML}{555555}
\definecolor{lightbg}{HTML}{F7F7F7}
\definecolor{borderc}{HTML}{E0E0E0}

% ---------------------------------------------------------
% Numbered circle
% ---------------------------------------------------------

\newcommand{\circnum}[2]{%
  \tikz[baseline=-0.65ex]{%
    \node[
      circle,
      fill=#1,
      text=white,
      font=\bfseries\small,
      inner sep=0pt,
      minimum size=0.42cm
    ] {#2};
  }%
}

% ---------------------------------------------------------
% Annotation number on screenshot
% ---------------------------------------------------------

\newcommand{\annot}[4]{%
  \node[
    circle,
    fill=#4,
    draw=white,
    line width=1pt,
    minimum size=0.44cm,
    inner sep=0pt,
    font=\bfseries\scriptsize,
    text=white
  ]
  at ($(img.north west)+(#1,-#2)$) {#3};
}

% ---------------------------------------------------------
% Legend item
% ---------------------------------------------------------

\newcommand{\legenditem}[4]{%
  \par\noindent
  \circnum{#1}{#2}
  \hspace{3pt}
  {\scriptsize\bfseries\color{darktext}#3}\\
  {\tiny\color{graytext}\RaggedRight #4\par}
  \vspace{1pt}
}

% ---------------------------------------------------------
% Reflection box
% ---------------------------------------------------------

\newcommand{\reflectionbox}[2]{%
  \begin{tcolorbox}[
    colback=lightbg,
    colframe=borderc,
    boxrule=0.5pt,
    arc=2pt,
    left=7pt,
    right=7pt,
    top=3pt,
    bottom=3pt
  ]
    \scriptsize
    \textbf{\color{darktext}#1}
    {\color{graytext}#2}
  \end{tcolorbox}
}

% =========================================================
% DOCUMENT
% =========================================================

\begin{document}

% =========================================================
% SLIDE 1 - HCI
% =========================================================

\begin{frame}[t]

\vspace*{-1mm}

{\large\bfseries\color{darktext}
Screenshot 1 --- HCI Focus}
\quad
{\small\color{graytext}
(Tinkercad, 3D Design Interface)}

\vspace{2mm}

\begin{columns}[t]

% ---------------------------------------------------------
% LEFT SIDE - SCREENSHOT
% ---------------------------------------------------------

\begin{column}{0.60\textwidth}

\begin{tikzpicture}

\node[
  anchor=north west,
  inner sep=0
] (img) at (0,0)
{\includegraphics[
  width=7.2cm
]{hci.png}};

\draw[
  borderc,
  line width=0.6pt
]
(img.north west)
rectangle
(img.south east);

% Annotation 1 - View Cube
\annot{0.30cm}{0.96cm}{1}{hciblue}

% Annotation 2 - Zoom controls
\annot{0.13cm}{2.01cm}{2}{hciblue}

% Annotation 3 - Shape Gallery
\annot{6.27cm}{1.48cm}{3}{hciblue}

\end{tikzpicture}

\end{column}

% ---------------------------------------------------------
% RIGHT SIDE - HCI INFORMATION
% ---------------------------------------------------------

\begin{column}{0.40\textwidth}

{\scriptsize\bfseries\color{hciblue}
HCI ELEMENTS}

\vspace{1mm}

\legenditem{hciblue}{1}
{Orientation (View) Cube}
{Click a face to snap the camera to a standard view --- a direct way to reorient.}

\legenditem{hciblue}{2}
{Zoom +/- Controls}
{Large, high-contrast buttons that clearly signal ``clickable'' for a common action.}

\legenditem{hciblue}{3}
{Shape Gallery + Search}
{Drag-and-drop panel with search, letting users recognize shapes instead of recalling names.}

\end{column}

\end{columns}

\vfill

% ---------------------------------------------------------
% REFLECTION 2
% ---------------------------------------------------------

\reflectionbox
{Reflection 2 --- HCI Improvement: }
{I would add clearer labels or tooltips to the main navigation controls.
This would help beginners understand the purpose of each control more quickly.}

\end{frame}


% =========================================================
% SLIDE 2 - COMPUTER GRAPHICS
% =========================================================

\begin{frame}[t]

\vspace*{-1mm}

{\large\bfseries\color{darktext}
Screenshot 2 --- CG Focus}
\quad
{\small\color{graytext}
(Tinkercad, 3D Design Interface)}

\vspace{2mm}

\begin{columns}[t]

% ---------------------------------------------------------
% LEFT SIDE - SCREENSHOT
% ---------------------------------------------------------

\begin{column}{0.60\textwidth}

\begin{tikzpicture}

\node[
  anchor=north west,
  inner sep=0
] (img) at (0,0)
{\includegraphics[
  width=7.2cm
]{cg.png}};

\draw[
  borderc,
  line width=0.6pt
]
(img.north west)
rectangle
(img.south east);

% Annotation 1 - 3D Meshes
\annot{2.92cm}{1.57cm}{1}{cgred}

% Annotation 2 - Shading
\annot{2.30cm}{1.97cm}{2}{cgred}

% Annotation 3 - Perspective Grid
\annot{4.75cm}{0.99cm}{3}{cgred}

\end{tikzpicture}

\end{column}

% ---------------------------------------------------------
% RIGHT SIDE - CG INFORMATION
% ---------------------------------------------------------

\begin{column}{0.40\textwidth}

{\scriptsize\bfseries\color{cgred}
CG ELEMENTS}

\vspace{1mm}

\legenditem{cgred}{1}
{3D Primitive Meshes}
{Each shape is a polygon mesh --- vertices and faces defining a 3D solid.}

\legenditem{cgred}{2}
{Shading \& Specular Highlight}
{The sphere's light-to-dark gradient and highlight simulate a light source, reading as ``round''.}

\legenditem{cgred}{3}
{Perspective Grid (Workplane)}
{Grid lines converge toward the horizon, signaling a perspective camera and depth.}

\end{column}

\end{columns}

\vfill

% ---------------------------------------------------------
% REFLECTION 1
% ---------------------------------------------------------

\reflectionbox
{Reflection 1 --- CG Quality: }
{Shading and perspective help shapes read as real 3D objects, so size and overlap are judged intuitively.
However, perspective distortion can hinder precision work because edge objects may appear skewed, complicating exact alignment.}

\end{frame}

\end{document}